\documentclass[conference]{IEEEtran}
\usepackage{newtxtext,newtxmath}
\usepackage{amsmath}
\usepackage{booktabs}
\usepackage{multirow}
\usepackage{array}
\usepackage{url}
\usepackage{cite}
\usepackage[hidelinks]{hyperref}
\newcommand{\st}[1]{\textsc{#1}}

\begin{document}

\title{Learning Where to Look: Multimodal Deep Learning\\with Verified Outputs for Digital Forensic Triage}

\author{
\IEEEauthorblockN{Daniel Cabrera}
\IEEEauthorblockA{\textit{Departamento de Ingeniería}\\
\textit{Universidad del Pacífico}\\
Lima, Perú\\
ds.cabreral@alum.up.edu.pe}
\and
\IEEEauthorblockN{Fabrizzio Marin}
\IEEEauthorblockA{\textit{Departamento de Ingeniería}\\
\textit{Universidad del Pacífico}\\
Lima, Perú\\
sf.marinc@alum.up.edu.pe}
\and
\IEEEauthorblockN{Luis Moran}
\IEEEauthorblockA{\textit{Departamento de Economía}\\
\textit{Universidad del Pacífico}\\
Lima, Perú\\
l.moranl@alum.up.edu.pe}
}

\maketitle

\begin{abstract}
A forensic examiner facing a seized phone or disk must decide which of tens of thousands of files to read first,
and any model that helps must never put an invented fact into a report. We train three deep-learning models on a
corpus of 72 public forensic cases (162{,}793 files) and connect them in a triage pipeline whose outputs are checked
by a deterministic verifier. \emph{Jev}, a byte-level CNN, identifies the file type of 512-byte fragments without
their header (0.811 accuracy, 17 types). \emph{Jev-MM}, a transformer that fuses CLIP image features, container
metadata and OCR text, predicts where an image came from without seeing its path (0.759 accuracy for a three-seed
ensemble vs.\ 0.692 for logistic regression and 0.662 for the best single modality). The \emph{investigation
planner}, a five-modality network trained on the answer keys of 16 solved cases, ranks the files of an unseen case
so that the agents read the right ones first; it can use a relevance prior from a local LLM that we fine-tune with
LoRA on 284 public forensic write-ups. Under a pre-registered nested leave-one-case-out protocol, the selected planner
reaches 0.813 AUROC and 0.292 average precision on unseen cases (random ranking: 0.5 and 0.032), and its 100
top-ranked files hold 75\,\% of a case's answer items. Three model-driven agents (planner, LLM analyst, image agent)
orchestrated around a blackboard of verified leads recover 136 of 869 answer-key identifiers reading 100 files per
case, about as many as deterministic parsers that read every file (137); the verifier rejects 545 of 1{,}707 LLM
quotes, so no hallucinated quote becomes evidence.
\end{abstract}

\begin{IEEEkeywords}
multimodal learning, digital forensics, learning to rank, self-supervised learning, calibration, hallucination
\end{IEEEkeywords}

\section{Introduction}

Digital investigations routinely involve several computers, phones and USB drives, each with tens of thousands of
files, thousands of SQLite databases and hundreds of thousands of file-system timestamps \cite{garfinkel2010,quick2014}.
Manual review has become the bottleneck of forensic laboratories. Deep learning offers two kinds of help: models that
\emph{recognize} content (file types, image provenance, text in screenshots) and models that \emph{prioritize} it,
telling the examiner where the answer to an investigation is likely to be. Large language and vision-language models
\cite{llava,qwen25vl} add fluent descriptions, but they may also produce claims that their input does not support
\cite{kalai2025}; in a forensic report a single invented timestamp can compromise the admissibility of the whole
analysis.

This paper treats forensics as the application domain of a deep-learning study. We ask three questions. (i) Can
small models trained on public forensic data recognize evidence when the cues examiners rely on (headers, paths)
are missing? (ii) Does multimodal fusion help, and when? (iii) Can a model learn, from solved cases, where to look in
a new one, and can its output be used without risking hallucinated evidence? Our contributions are:
\begin{itemize}
\item \textbf{Datasets without manual labelling.} From a corpus of 108 registered public cases (72 ready, 162{,}793
extracted files) we derive four datasets whose labels come from deterministic rules or published answer keys: 145{,}589
file fragments, 12{,}532 images, 33{,}047 files with relevance labels from 16 solved cases, and 80{,}126 unlabelled
files for self-supervised pretraining. All splits are by case.
\item \textbf{Jev}, a 0.89\,M-parameter byte CNN with typed, calibrated heads (file type, disguise, position).
\item \textbf{Jev-MM}, a multimodal transformer over pixels, metadata and OCR text that predicts image provenance
without access to the file path, with a full modality ablation over three seeds.
\item \textbf{An investigation planner} that ranks files by predicted relevance from five modalities, and a
controlled study of multimodal fusion under severe case scarcity: plain fusion, self-supervised contrastive
pretraining, metadata anchoring, wide-and-deep residual learning, data augmentation with synthetic positives,
transfer learning from Jev, and noisy-student self-training.
\item \textbf{A verified multi-agent pipeline}: the planner orchestrates an LLM analyst and an image agent in rounds;
every statement carries machine-checkable citations and a deterministic engine re-reads the original evidence before
anything is reported as evidence. Model outputs without citations are kept as labelled leads.
\item \textbf{A pre-registered evaluation} reported as one benchmark table per task, including the negative results.
\end{itemize}

\section{Background}

\textbf{Convolutional and transformer encoders.} One-dimensional CNNs over byte or character embeddings capture
local patterns such as magic numbers, record structures and path components; global max and mean pooling make them
length-independent. Transformers relate tokens through attention \cite{vaswani2017}; in multimodal fusion each
modality is encoded into one or more tokens and a transformer encoder with a learned classification token
(\texttt{CLS}) mixes them, while a key-padding mask removes absent modalities.

\textbf{Pretrained encoders.} CLIP \cite{clip} learns image and text encoders by contrastive alignment of 400\,M
image--caption pairs; its vision transformer gives transferable image features. Sentence encoders such as
multilingual MPNet \cite{sbert} map text to fixed vectors. Contrastive objectives (InfoNCE \cite{oord2018}) can also
align different \emph{views of the same object} without labels, which we use for self-supervised pretraining.

\textbf{Calibration.} Neural classifiers are often over-confident; temperature scaling divides the logits by a
scalar fitted on validation data and lowers the expected calibration error (ECE) \cite{guo2017}. Calibrated
probabilities allow a model to abstain.

\textbf{Multimodal imbalance.} A multimodal network can be worse than its best single-modality counterpart: the
modalities learn at different speeds and the network overfits the easiest ones \cite{wang2020}. Remedies include
per-modality losses, modality dropout and architectures that keep a strong unimodal path.

\textbf{Hallucination and grounding.} A hallucination is generated content that is not supported by the input
\cite{kalai2025,ji2023}. We ground model outputs by requiring each reported fact to cite a location in the original
evidence that a deterministic program can re-read; cryptographic hashes (SHA-256) show that the cited file has not
changed \cite{carrier2005}.

\section{Related Work}

\textbf{LLMs and agents in forensics.} Wickramasekara and Scanlon \cite{wickramasekara2024} propose four AutoGen agents
backed by LLaMA and StarCoder to automate forensic phases, without an empirical evaluation; they list accuracy and
hallucinations as open problems. Khatiwala \emph{et al.} \cite{khatiwala2026} use GPT-4 with a forensic knowledge
graph and SHA-256 identifiers on one Samsung phone image, reaching 95.24\,\% artifact precision, but depend on a cloud
model. ForensicLLM \cite{forensicllm} fine-tunes a 4-bit LLaMA-3.1-8B with retrieval-augmented fine-tuning and
attributes sources correctly 86.6\,\% of the time; it answers questions about forensics rather than analysing case
evidence. In contrast, our models are small, trained on case evidence, run locally, and their outputs are verified
against the evidence at the level of content, not only file integrity.

\textbf{Deep learning on forensic data.} File-fragment classification with CNNs over raw bytes is established
\cite{mittal2021}; Jev follows this line and adds calibrated typed heads and a disguise detector. Guwor and Shabaz
\cite{guwor2026} fuse MobileNetV2 visual features with EXIF-based features for forgery detection (83.85\,\% vs.\
68.61\,\% visual-only and 75.85\,\% metadata-only), showing that pixels and metadata are complementary. Jev-MM fuses
the same two modalities plus OCR text, with a transformer instead of stacked classifiers, for a different question
(provenance rather than authenticity) and with case-level splits.

\textbf{Learning with few labelled groups.} When labelled data are scarce, wide-and-deep models \cite{cheng2016}
combine a linear model with a deep network trained jointly; residual gates initialised at zero (ReZero
\cite{bachlechner2021}) start a network as the identity of its strongest path. Data augmentation and synthetic
minority examples (SMOTE \cite{chawla2002}, mixup \cite{zhang2018}) enlarge the training set; noisy-student
self-training \cite{xie2020} turns unlabelled data into soft labels. We evaluate all of these on the planner, where
the independent unit is a case and only 16 are labelled.

\section{Data}

\subsection{Forensic Corpus}
The corpus registers 108 public cases from NIST CFReDS, Digital Corpora, the DFRWS challenges, CTF competitions
(Magnet, Magnet Virtual Summit, Cellebrite, NFI, DFIR Madness), the Ali Hadi challenges, the VTO drone program and
the Hickman reference phones. A case groups the evidence items of one scenario, the brief given to examiners and,
when published, the answer key. 72 cases are ready (at least one extracted item and a usable brief): 100 evidence
items, 162{,}793 files (43.1\,GB), 2{,}813{,}350 file-system entries and 7{,}373 recovered deleted files, covering
Android (40\,\% of files), Windows (25\,\%), iOS (17\,\%), Linux (10\,\%), removable media, drones, cloud exports,
memory and network captures (Table~\ref{tab:data}). Disk images (E01, split RAW, LVM2, XFS, APFS) were opened with The
Sleuth Kit and \texttt{dissect}; large archives were read by HTTP range requests, so only selected members were
downloaded and each was checked by CRC-32. Every file's SHA-256 is recorded. Agents and models receive only the brief;
answer keys are stored apart and used only for labels and scoring.

\textbf{Integrity audit.} Because the corpus lived in a cloud-synchronized folder, we re-hashed all 162{,}793 files.
51 had changed or disappeared (29 \texttt{desktop.ini} files rewritten by Windows, 17 SQLite \texttt{-shm} files
rewritten, 5 missing). 39 were re-acquired from the original images with matching hashes, 22{,}395 at-risk files
were made read-only, and 12 that sit in images compressed at the source were excluded. The same audit found that the
fragment dataset builder had silently skipped 6{,}036 files with paths longer than 260 characters; the dataset was
rebuilt and Jev retrained.

\begin{table}[t]
\caption{Datasets Derived from the Corpus (All Splits by Case)}
\label{tab:data}
\centering
\small
\setlength{\tabcolsep}{4pt}
\begin{tabular}{@{}p{1.75cm}p{3.75cm}r@{}}
\toprule
\textbf{Dataset} & \textbf{Label source} & \textbf{Size} \\
\midrule
Fragments (Jev) & header signature of the file; extension rule; offset & 145{,}589 frag. \\
Images (Jev-MM) & storage location (path rules); path never an input & 12{,}532 img. \\
Planner & answer key of 16 solved cases (file named or identifier in bytes; 849 relevant) & 33{,}047 files \\
Pretraining & none (50 cases without key) & 80{,}126 files \\
\bottomrule
\end{tabular}
\end{table}

\subsection{Derived Datasets}
\textbf{Fragments.} 512-byte, sector-aligned fragments from file \emph{bodies} (the first sector, which carries the
signature that produced the label, is never used): 145{,}589 fragments from 46{,}557 distinct files, deduplicated by
SHA-256 across splits. Labels: 17 file types; whether the file is disguised (content contradicts the extension);
relative position in the file. Extensions are re-drawn from a shared pool during training so that the extension alone
cannot reveal the type.

\textbf{Images.} 12{,}532 decodable images from 57 cases (train / validation / test 9{,}914 / 808 / 1{,}810; at most
400 per case and class). The label is the origin implied by where the device stored the image: camera
(\texttt{DCIM}), screenshot, messaging media, user file, web cache, thumbnail or application asset; a second label
says whether the image is user-generated. The path is never a model input, so the model must infer provenance from
content, as for a carved or renamed image. 4{,}293 images carry OCR text.

\textbf{Planner.} A file of a solved case is \emph{relevant} if the answer key names it or its bytes contain one of
the key's identifiers (e-mail address, URL, domain, phone number, file name), found with an Aho--Corasick scan over
ASCII and UTF-16 encodings. For two cases whose keys are narrative (MVS~2024, NGDC~2012) we transcribed distinctive
solution tokens by hand, with their source. Result: 16 cases, 33{,}047 files, 849 relevant (2.6\,\%). Each held-out
case also keeps its set of \emph{answer items} (identifiers and named files) for coverage metrics.

\textbf{Pretraining.} The 50 ready cases without an answer key provide 80{,}126 unlabelled files after SHA-256
deduplication; no labelled case is included.

\section{Methods}

\subsection{Jev: Byte-Level Fragment Model}\label{sec:jev}
Each byte is embedded ($256\!\to\!64$) and passed through four Conv1d--BatchNorm--GELU blocks (kernels 9, 7, 5, 3,
with max pooling) and mean$\|$max pooling into $z\in\mathbb{R}^{256}$ (0.89\,M parameters). Three typed heads read
$z$: \emph{choice} (17 types), \emph{yes/no} (disguised; also reads a character embedding of the extension and the
predicted type distribution) and \emph{score} (position, five ordinal bins). The loss is
$\mathcal{L}=\mathrm{CE}_{\text{ls}=0.05}+0.5\,\mathrm{BCE}+0.3\,\mathrm{CE}_{\text{ord}}$; training uses AdamW, a
one-cycle schedule, bf16 and class-balanced sampling ($w\propto n^{-1/2}$) for 20 epochs, with the epoch chosen on
validation and a temperature fitted per head. Answers below 0.5 confidence abstain.

\subsection{Jev-MM: Multimodal Image Provenance}\label{sec:jevmm}
For an image $x$ the model receives three inputs taken from the file itself:
\begin{itemize}
\item \emph{Vision}: a $224\times224$ letterboxed thumbnail encoded by CLIP ViT-B/32 \cite{clip}; only the last
transformer block is fine-tuned.
\item \emph{Metadata}: 43 container features: format, geometry, size and bits per pixel, color mode, the EXIF tags
present (make, model, GPS, capture time, software, UTC offset, exposure), the JPEG quality implied by the
quantization table, chroma subsampling, ICC profile and PNG text chunks.
\item \emph{Text}: EasyOCR text embedded by a frozen multilingual sentence encoder \cite{sbert}, plus nine
statistics (boxes, confidence, text area, presence of a status-bar clock).
\end{itemize}
Each modality is projected to a token $h_m\in\mathbb{R}^{256}$; a learned $h_{\text{cls}}$ and the modality tokens
(plus modality embeddings $e_m$) pass through a two-layer transformer encoder in which absent modalities are masked:
\begin{equation}
z = \mathrm{Enc}\big([h_{\text{cls}};\, h_v+e_v;\, h_m+e_m;\, h_t+e_t],\ \text{mask}\big)_{\text{cls}},
\end{equation}
\begin{equation}
p(\text{orig.}\mid x)=\mathrm{softmax}(W_o z / T_o),\ \ p(\text{user}\mid x)=\sigma(w_u^\top z / T_u).
\end{equation}
The loss is label-smoothed cross-entropy plus $0.5\times$ BCE; each modality is dropped with probability 0.25 during
training so the model answers with whatever is available. Training: AdamW, one-cycle, bf16, class-balanced sampling,
12 epochs (9.8\,M trainable of 90.2\,M parameters). The number of fine-tuned ViT blocks (0/1/2) and its learning rate
were chosen on validation only. Temperatures $T_o,T_u$ are fitted on validation \cite{guo2017}.

\subsection{Investigation Planner}\label{sec:planner}
For each file $f$ of a case $C$ the planner predicts $p(\text{relevant}\mid f, C)$ from five modalities, each encoded
into one 192-dimensional token:
\begin{itemize}
\item \emph{path}: the last 160 characters, right-aligned so the file name is always at the end (character CNN);
\item \emph{header}: the first 256 bytes (byte CNN);
\item \emph{metadata}: a 47-dimensional vector: log size, depth, category, platform, deleted and hidden flags, MACB
timestamps as within-case percentiles (MLP);
\item \emph{content}: printable strings of the first 64\,KB embedded by the sentence encoder (MLP);
\item \emph{context}: the case brief, same encoder (MLP).
\end{itemize}
A two-layer transformer (four heads) fuses $[\text{CLS};\text{tokens}]$ and a linear head gives the logit. Training
uses binary cross-entropy with case-balanced sampling (each case contributes equally, half of its draws positive),
modality dropout 0.15, AdamW and a one-cycle schedule for 10 epochs of 150 steps.

Because the independent unit is the case and only 16 are labelled, plain fusion is at risk of fitting case-specific
cues (one case's URLs or folder names). We study the following remedies, all evaluated under the same protocol.

\textbf{Self-supervised encoders.} The four file encoders are pretrained on the 80{,}126 unlabelled files with a
symmetric InfoNCE loss over the six modality pairs: the views of one file (path, header, metadata, content) must
identify each other among the 512 files of a batch, as in CLIP but between modalities of one file. On six held-out
unlabelled cases the encoders retrieve the matching view with 10--35\,\% top-1 accuracy (chance 0.2\,\%). The
planner loads them frozen or fine-tuned at $0.1\times$ the learning rate.

\textbf{Metadata anchoring.} The logit is $a(t_m)+g\,h(z_{\text{cls}})$, where $t_m$ is the metadata token before
fusion, $a$ and $h$ are linear heads and $g$ is a scalar gate initialised at zero \cite{bachlechner2021}; an auxiliary
BCE trains $a$ alone. The network starts as a metadata model and fusion can only add a correction.

\textbf{Wide and deep.} The anchor is replaced by the linear model itself \cite{cheng2016}: the logit is
$w^\top m + b + g\,h(z_{\text{cls}})$, where $(w,b)$ is the L2-regularized logistic regression on the raw metadata
vector $m$, fitted on the fold's non-test cases and frozen, and the deep fusion learns a gated residual on the training
cases. Epoch~0 ($g=0$, the linear model alone) is a candidate in the validation-based epoch choice, so the deep part is
kept only if it improves validation ranking.

\textbf{Augmentation and synthetic positives.} Per batch: random prefix crop of the path and 3\,\% character
substitutions outside the file name; 2\,\% byte noise after the 16-byte signature; Gaussian jitter on the continuous
metadata columns; dropout and noise on the content embedding; and SMOTE-style synthetic positives \cite{chawla2002}:
half of the positives get metadata and content interpolated towards another positive,
$\lambda\sim\mathcal{U}(0.5,1)$.

\textbf{Transfer learning.} The header encoder is replaced by the convolutional trunk of Jev (trained on 145k
fragments) and fine-tuned at $0.1\times$ the learning rate.

\textbf{Noisy-student self-training.} An anchored teacher trained on the labelled training cases scores the 80{,}126
unlabelled files; the student is trained on the labelled batch plus a half batch of unlabelled files with the
teacher's soft labels (weight 0.5) under augmentation \cite{xie2020}.

\textbf{LLM relevance prior.} A local LLM (Qwen2.5-VL-3B, text only \cite{qwen25vl}) sees the case brief and, for one
file, its path, category, size and a 500-character text excerpt, and is asked whether the file is likely to answer the
investigation. Its prior is $L=\log P(\text{Yes})-\log P(\text{No})$ from one forward pass (fp32 logits of the two
answer tokens, no generation). In the wide part, $m$ is extended with $L/10$ and $L$ standardised within the case (no
labels are used). The prior comes in two versions with the same prompt. \emph{Zero-shot}: the LLM as released.
\emph{Fine-tuned}: the same LLM after LoRA fine-tuning \cite{hu2022lora} on public write-ups of other forensic
challenges, so that it learns from examiners' own accounts where they looked and which tools they used. We collected
284 English Markdown write-ups (CTF and training labs: Windows triage images, memory, logs, e-mail, mobile) and kept each
as one plain-text pair, \emph{case context} (the scenario before the first question) $\rightarrow$ \emph{solution}
(the rest), with no annotation. The solution text was cleaned so that the examiner's actions dominate it: it is laid out
as question $\rightarrow$ the examiner's own steps (which artifact was opened, with which tool, what was searched)
$\rightarrow$ answer; questions answered without any step, links, images, page metadata, summaries and pasted tool
output beyond eight lines were removed. A write-up was discarded if its raw page named any of our 16 labelled cases or
contained any of their answer-key identifiers (108 discarded), or was not in English. Fine-tuning used a recipe fixed in advance: rank 16, $\alpha=32$,
dropout 0.05 on all attention and MLP projections of the language model (30M trainable parameters), learning rate
$10^{-4}$, one epoch, 1{,}024 tokens per example, loss on the solution tokens only. Since no labelled case is seen, one
fine-tuned model serves every fold.

\subsection{Verified Triage Pipeline}\label{sec:pipeline}
Given a case, the planner ranks all files and the top $K$ are dispatched: images to an image-triage agent (Jev-MM)
and, optionally, a locally run vision-language model (Qwen2.5-VL-3B in 4-bit NF4 \cite{qwen25vl,qlora}); other files
to a reader agent that reports e-mail addresses and URLs. Every statement is a \emph{finding} with a set of citations
$C_f$; each citation gives the file, its expected SHA-256 $h_c$, a locator (byte offset, line, read-only SQL query,
regular expression, EXIF tag, registry value or event-log record) and the expected snippet. A deterministic engine,
with its own file access and hash cache, recomputes the hash, runs the locator on the original file and compares the
result. Its status is
\begin{equation}
\mathrm{st}(f) =
\begin{cases}
\st{Uncorroborated}, & C_f = \emptyset,\\
\st{Verified}, & \bigwedge_{c \in C_f} v(c) \wedge \kappa(f),\\
\st{Rejected}, & \text{otherwise},
\end{cases}
\label{eq:status}
\end{equation}
where $v(c)$ is the outcome of re-running citation $c$. For text read by a model (OCR or VLM), $\kappa$ requires a
fuzzy match against a database cell or text line of the case: with $M(a,b)$ the matching-block length of
Ratcliff--Obershelp \cite{ratcliff1988}, $\kappa = (|a|\ge 10) \wedge (M/|a| \ge 0.85) \wedge (|a|\ge 25 \vee M/|b|\ge
0.5)$; otherwise $\kappa=1$. Only \st{Verified} findings enter the evidence sections of a report. Outputs of Jev and
Jev-MM carry no citations and are therefore always \st{Uncorroborated} leads: the models prioritize, the verifier
decides. The evidence is never written; SQLite files are queried on hash-keyed copies.

\subsection{Orchestrated Multi-Agent Investigation}\label{sec:orch}
The pipeline above reads the top $K$ files once. The orchestrated system replaces it with three model-driven agents
that work in rounds around a shared \emph{blackboard} of verified leads (Fig.~\ref{fig:orch}).

\textbf{Orchestrator.} The planner ranks every file of the case. In round $r$ it dispatches the next $B$ files (here
$B=25$, $K=100$, so four rounds): non-image files to the analyst agent, images to the image agent. From round 2 on,
files are re-ranked with the leads $L$ the agents have verified so far:
\begin{equation}
s_i = \operatorname{logit}(p_i) + \beta \log\bigl(1 + |\{\ell \in L : \ell \text{ occurs in file } i\}|\bigr),
\label{eq:rerank}
\end{equation}
where occurrence is tested on the text snippet of the file's index entry and $\beta=2$ is fixed a priori (never tuned
on a case). Read files never return.

\textbf{Analyst agent.} A local instruction-following LLM (Qwen2.5-VL-3B used text-only \cite{qwen25vl}) reads up to
3{,}500 characters of each file, taking windows around current leads first, together with the case brief and the
leads on the blackboard. It answers in JSON with up to six \emph{verbatim} quotes (typed as person, account, message,
location, \ldots) and new leads. Each quote is searched in the file's bytes (UTF-8, UTF-16LE, then case-insensitive);
when found, the finding cites the byte offset and the claim is written from the file's own bytes, never from the
model's text. When it is not found, the finding cites the model's quote anyway, so the engine \st{Rejects} it: a
hallucination caught and counted. A lead enters the blackboard only if it occurs inside a verified quote. The analyst
also calls a deterministic identifier-extraction tool (e-mail addresses, URLs).

\textbf{Image agent.} Jev-MM triages every dispatched image; the vision-language model reads the three images per
round most likely to be user-generated. Text read from an image becomes evidence and a lead only when the engine
corroborates it against a case artifact ($\kappa$ above).

\begin{figure}[t]
\centering
\small
\setlength{\fboxsep}{3pt}
\fbox{\parbox{0.92\columnwidth}{\centering
Planner (orchestrator): rank all files $\rightarrow$ pick next $B$ by Eq.~\eqref{eq:rerank}\\[2pt]
$\downarrow$ text files \hspace{6em} $\downarrow$ images\\[2pt]
Analyst LLM (quotes, leads) \hspace{1em} Jev-MM $\rightarrow$ VLM (top 3)\\[2pt]
$\downarrow$ \hspace{10em} $\downarrow$\\[2pt]
Verification engine (re-reads evidence bytes)\\[2pt]
$\downarrow$ \textsc{Verified} only\\[2pt]
Blackboard of leads $\rightarrow$ next round}}
\caption{Orchestrated investigation. Only verified statements reach the blackboard or the report.}
\label{fig:orch}
\end{figure}

\section{Experimental Setup}
\textbf{Jev and Jev-MM.} Fixed case-level splits shared by both models (4 validation and 4 test cases: two Android and two iOS reference phones, two Windows scenarios and two CTF images, so every platform family is represented). Jev-MM and
each modality subset are trained with seeds 13, 14 and 15 and reported as mean $\pm$ std on the test cases; the
deployed checkpoint is additionally evaluated with modalities withheld at inference. Baselines: majority class,
logistic regression on the byte histogram (Jev), extension-only and rule-based disguise detectors.

\textbf{Planner (pre-registered).} The protocol and the candidates were written down and committed before any of
the numbers below were computed; one later addition (the fine-tuned prior) was registered before any
held-out number was read. Earlier exploratory runs had compared many variants on the held-out cases themselves; they
informed the candidate list and are not reported as evidence. Leave-one-case-out (LOCO): each labelled case $t$ is
held out once and the other 15 form its pool. \emph{Nested selection}: every pool case with at least 10 relevant
files serves in turn as inner-validation case $v$; each candidate (metadata only, anchored, wide and deep, wide and
deep with augmentation, and wide and deep with the zero-shot or the fine-tuned LLM prior) is trained on pool$\,-\,v$
with seeds 13, 14 and 15 for 10 epochs, and the linear part of wide and deep is fitted on pool$\,-\,v$ only. The
(candidate, epoch) with the highest mean inner-validation AP is selected, and the final model of fold $t$ is the
mean-logit ensemble of all its inner models at that epoch. This one final model per fold is the planner we report
and deploy; per-candidate numbers are descriptive. Metrics are computed on the 13 held-out cases where ranking is
meaningful (at least 50 files and not all relevant), each case counting once: average precision (AP, primary), AUROC
and coverage@$K$, the share of the case's answer items in its $K$ top-ranked files. A random ranking has AUROC 0.5
and expected AP equal to the share of relevant files (0.032 on average over the 13 cases). \emph{Secondary metric}:
the same scores restricted to user-level files, i.e., excluding containers (file-system metadata, memory images,
network captures, raw disk and NAND images, paging files and unallocated space). The primary labels favour
containers: a container or a large database contains answer-key bytes almost by construction, and its size and
category are enough to rank it high.

\textbf{End to end (pre-registered).} Each solved case is investigated by the planner fold that never saw it, reading
at most $K=100$ files ($K$ was fixed from the timing of one case, before any comparison). The metric is the number of
answer-key identifiers that appear in \st{Verified} findings, counted separately by source: verbatim LLM quotes, the
deterministic extraction tool, the image agent, and their union. Cases without automatic answer-key identifiers are
excluded: MVS 2024 and NGDC 2012 (labelled only through curated tokens) and DFRWS 2017 (no identifiers). The system
uses $K=100$, $B=25$, the unquantised LLM and a GPU used by nothing else. One run stopped on a GPU out-of-memory error
and was resumed with the same configuration (allocator setting only). As a reference, three deterministic parsers
(application databases, file system, Windows artifacts) read every file of every case.

\textbf{Hardware.} One laptop with an NVIDIA RTX 5080 Laptop GPU (16\,GB, capped at 12\,GB in software) and 32\,GB
RAM; no data leave the machine.

\section{Results}

\subsection{Jev}
On the test cases Jev identifies the type of a headerless 512-byte fragment with 0.811 accuracy and 0.658 macro-F1,
against 0.630 / 0.470 for logistic regression on the byte histogram and 0.241 for the majority class. Calibrated, it
answers 91\,\% of fragments with 85\,\% accuracy. Structured formats are recognized well (event logs F1 0.94, text
0.95, registry hives 0.93, LNK 0.86, JPEG 0.84); compressed bodies (MP4, ZIP, GIF) are near-random bytes and remain
hard (F1 0.03--0.38). The disguise head reaches 0.979 AUROC with re-drawn extensions, but 0.758 on real extensions,
where disguised files are rare (103 of 9{,}591; precision about 5\,\% at 0.5). An early version reached 0.97 AUROC
from the extension alone; the shared extension pool removed that shortcut, and the extension-only baseline remains in
the evaluation as a check.

\begin{table}[t]
\caption{Image Provenance Benchmark on the Held-Out Test Cases (Mean $\pm$ Std over 3 Seeds; LR = Logistic Regression)}
\label{tab:jevmm}
\centering
\small
\setlength{\tabcolsep}{3pt}
\begin{tabular}{@{}lccc@{}}
\toprule
\textbf{Model} & \textbf{Accuracy} & \textbf{Macro-F1} & \textbf{User AUC} \\
\midrule
Vision only & $0.662 \pm 0.010$ & $0.423 \pm 0.012$ & $0.850 \pm 0.014$ \\
Metadata only & $0.655 \pm 0.035$ & $0.480 \pm 0.056$ & $\mathbf{0.909 \pm 0.025}$ \\
Text only & $0.502 \pm 0.017$ & $0.225 \pm 0.021$ & $0.611 \pm 0.009$ \\
Vision + metadata & $0.724 \pm 0.006$ & $0.520 \pm 0.007$ & $0.878 \pm 0.005$ \\
All three (Jev-MM) & $0.744 \pm 0.015$ & $0.533 \pm 0.015$ & $0.899 \pm 0.003$ \\
\textbf{Ensemble (3 seeds)} & $\mathbf{0.759}$ & $\mathbf{0.553}$ & $0.911$ \\
\midrule
LR, metadata & $0.655$ & $0.488$ & $\mathbf{0.915}$ \\
LR, metadata + text & $0.692$ & $0.489$ & $0.893$ \\
Majority class & $0.450$ & -- & $0.500$ \\
\midrule
\multicolumn{4}{@{}l}{\emph{One Jev-MM model, modalities withheld at inference}} \\
None withheld & 0.733 & 0.519 & 0.900 \\
Text & 0.735 & 0.543 & 0.888 \\
Metadata (no EXIF) & 0.647 & 0.396 & 0.847 \\
Vision & 0.697 & 0.515 & 0.916 \\
\bottomrule
\end{tabular}
\end{table}

\subsection{Jev-MM}
Table~\ref{tab:jevmm} shows that fusing the three modalities gives the best provenance classifier, 8 accuracy points
above the best single modality, and the most stable one across seeds (metadata alone varies by $\pm0.056$ macro-F1).
Vision and metadata are complementary: EXIF, JPEG tables and sizes identify camera photos and thumbnails, pixels
identify application assets and web images, and the attention of the fusion token follows this (0.60 on metadata for
camera photos, 0.76 on vision for application assets). Text adds little on average because most images carry none,
and withholding it costs nothing, which is the intended effect of modality dropout; stripping EXIF costs 9 points.
For the binary user-generated question, metadata alone is as good as fusion, and logistic regression on metadata is
the best model (AUROC 0.915). Temperature scaling lowers the ECE from 0.119 to 0.065; the single model answers 90\,\%
of test images with 78\,\% accuracy. Against logistic regression, chosen with its L2 penalty on validation, deep
fusion clearly wins on the 7-class task: $+0.067$ accuracy and $+0.064$ macro-F1 for the ensemble over the best
linear model. Three training recipes fixed in advance (original; full class balance; 20 epochs) were compared on
validation by the seed ensemble; the original recipe won, and its 3-seed ensemble (mean logits) was tested once
and deployed. Per-class F1 of the ensemble: camera 0.87, application asset 0.88, thumbnail 0.78, web cache 0.66,
screenshot 0.34, user file 0.19, messaging 0.15. The weak classes are those whose label is not visible in the image
(a Downloads folder) or that have few training examples (106 messaging images).

\subsection{Investigation Planner (Pre-Registered)}
\textbf{Final model.} On the 13 rankable held-out cases the pre-registered final model reaches 0.813 AUROC and 0.292
AP (Table~\ref{tab:planner}), nine times the AP of a random ranking (0.032). Its 10, 50 and 100 top-ranked files hold
41\,\%, 65\,\% and 75\,\% of a case's answer items, while 100 files are on average 4\,\% of a case. Per-case AP varies
widely (from 0.45 on M57 Patents and 0.42 on NIST Data Leakage to near zero on cases with one or two relevant files),
which is expected from 13 independent cases. On the secondary metric (user-level files only, 11 rankable cases) it
reaches 0.803 AUROC and 0.213 AP.

\textbf{What selection chose.} In 12 of 13 folds the nested selection chose wide and deep with the zero-shot LLM prior,
in one (DFRWS 2018) the fine-tuned prior, and always at epoch 0, i.e., the wide part plus the LLM prior with the deep
residual at its initial value: on inner validation, the LLM prior was worth more than further training of the fusion
network.

\textbf{Benchmark.} The other rows of Table~\ref{tab:planner} give each candidate at the epoch with its own best
inner-validation AP. They are selected without looking at the held-out case but are descriptive only. The metadata
MLP has the highest held-out AP (0.350) and anchored fusion the best coverage@10 (0.537); the LLM prior gives the
highest AUROC (0.815). Inner validation preferred the prior, so the deployed planner optimises global ordering more
than the very top of the list.

\textbf{Fine-tuned prior.} LoRA fine-tuning on 256 write-ups lowered the loss on 28 held-out write-ups from 2.531 to
2.320 (perplexity 12.57 $\rightarrow$ 10.18): the LLM learned the language of examiners' solutions. As a relevance
prior it behaves like the zero-shot one (AP 0.294 vs.\ 0.291, AUROC 0.806 vs.\ 0.815).

\begin{table}[t]
\caption{Evidence Ranking Benchmark: 13 Held-Out Cases (LOCO)}
\label{tab:planner}
\centering
\small
\setlength{\tabcolsep}{3.5pt}
\begin{tabular}{@{}lccccc@{}}
\toprule
\textbf{Model} & \textbf{AUROC} & \textbf{AP} & \textbf{c@10} & \textbf{c@50} & \textbf{c@100} \\
\midrule
\textbf{Final (nested selection)} & \textbf{0.813} & 0.292 & 0.409 & 0.652 & 0.747 \\
Random ranking & 0.500 & 0.032 & -- & -- & -- \\
\midrule
\multicolumn{6}{@{}l}{\emph{Descriptive: each candidate at its inner-val.\ best epoch}} \\
Metadata MLP & 0.767 & \textbf{0.350} & 0.496 & 0.686 & 0.746 \\
Anchored & 0.764 & 0.326 & \textbf{0.537} & \textbf{0.699} & 0.737 \\
Wide \& deep & 0.756 & 0.298 & 0.440 & 0.697 & \textbf{0.791} \\
\ + augmentation & 0.766 & 0.301 & 0.487 & 0.660 & 0.755 \\
\ + LLM prior & 0.815 & 0.291 & 0.420 & 0.652 & 0.747 \\
\ + tuned prior & 0.806 & 0.294 & 0.450 & 0.663 & 0.724 \\
\midrule
\multicolumn{6}{@{}l}{\emph{Secondary: user-level files only (11 cases)}} \\
Final (nested selection) & 0.803 & 0.213 & 0.305 & 0.602 & 0.648 \\
\bottomrule
\end{tabular}
\end{table}

\subsection{Orchestrated Agents End to End}
On the 13 cases with automatic identifiers (869 identifiers), reading 100 files per case, the verified findings of
the orchestrated system contain 136 identifiers (15.7\,\%) and at least one in 7 cases (Table~\ref{tab:e2e}); three
deterministic parsers that read \emph{every} file of every case found 137 in 4 cases. Most identifiers enter through
the analyst's extraction tool (132); verified LLM quotes contribute 21, among them typed facts (persons, accounts,
messages) that the tool does not extract. The image agent contributes no identifier, as expected: identifiers are
text, and its role is triage. The verifier rejected 545 of the 1{,}707 quotes of the analyst (32\,\%): these are
hallucinated or paraphrased quotes, and none of them became evidence. The final system, run with the improved Jev-MM
ensemble, gives the same numbers case by case. In a vision proof of concept on a Pixel~3 (58 images), GPS
coordinates were re-checked against EXIF and a screenshot reading ``Super Secret Message'' was tied to a launcher
database row through a read-only SQL citation (40 verified, 0 rejected, 244 uncorroborated descriptions; peak GPU
memory 6.5\,GB).

\begin{table}[t]
\caption{End to End: Answer-Key Identifiers in Verified Findings (13 Cases, 869 Identifiers)}
\label{tab:e2e}
\centering
\small
\begin{tabular}{@{}lrr@{}}
\toprule
& \textbf{Agents} & \textbf{Parsers} \\
& \textbf{(top 100 files)} & \textbf{(all files)} \\
\midrule
Identifiers, verified LLM quotes & 21 & -- \\
Identifiers, extraction tool & 132 & -- \\
Identifiers, image agent & 0 & -- \\
Identifiers, union & 136 & 137 \\
Cases with at least one & \textbf{7} & 4 \\
Files read per case & $\le$ 100 & all \\
LLM quotes rejected by the verifier & 545 / 1{,}707 & -- \\
\bottomrule
\end{tabular}
\end{table}

\section{Discussion}
The two multimodal models give different answers to the question ``does fusion help?'', and the difference is
instructive. For images, thousands of independent examples from 57 cases let the transformer learn complementary
cues, and fusion wins directly. For the planner, the independent unit is the case: with 15 training cases per fold,
the high-capacity path and content encoders memorise case-specific cues that do not transfer, the modality-imbalance
failure described by Wang \emph{et al.} \cite{wang2020}. Remedies that \emph{protect the strongest modality}
(anchoring, wide and deep) helped more than those that only add data or regularise the encoders, and under honest
selection on inner validation the strongest addition was not more fusion training but the LLM prior, which improves
the global ordering of files. Fine-tuning the LLM on write-ups made it a better model of examiners' text but not a
better relevance prior. With 16 labelled cases, per-case AP varies more than the differences between candidates, so
more labelled cases are the most direct way to improve the planner.

End to end, the agents read at most 100 files per case and reach as many identifiers as parsers that read every
file, in more cases. What the verifier adds is safety: about three in ten LLM quotes were not in the file, and each one
was rejected.

The verifier is what makes these models usable in a forensic setting. None of the trained models writes evidence:
Jev and Jev-MM produce calibrated, labelled leads, the planner only decides reading order, and every statement that
reaches the evidence sections was re-read from the original file by a program that shares nothing with the model. A
model error cannot silently become evidence.

\section{Limitations}
First, the labels are proxies: storage location for Jev-MM and automatically extracted answer-key identifiers for
the planner and the end-to-end metric, not examiner judgments; a file that the key does not mention may still be
relevant. Second, only 16 cases have usable answer keys, so planner estimates have wide confidence intervals and the
LOCO protocol is the only honest option. Third, verification shows that a cited value exists in the evidence, not
that its interpretation is correct. Fourth, the extraction tool covers only e-mail addresses and URLs; memory images,
network captures and drone flight logs have no model or parser. The base LLM may have seen write-ups of the dataset
cases during pre-training, which we cannot control. Fifth, GPU attention kernels are not fully deterministic, so
single runs vary; we report three seeds throughout. Finally, the corpus integrity audit showed that
cloud-synchronized storage alters evidence; court-style use requires a local, non-synchronized evidence store.

\section{Conclusion}
We trained three deep-learning models on public forensic cases and combined them in a multi-agent pipeline whose
outputs are verified against the original evidence. Where examples are plentiful, deep learning wins clearly: a byte
CNN recognizes headerless fragments (0.811 vs.\ 0.630 accuracy for a histogram baseline) and a multimodal transformer
recognizes image provenance better than any single modality (0.759 accuracy). The investigation planner, learned from
16 solved cases, ranks an unseen case so that its 100 top files hold three quarters of the answer items, and the
orchestrated agents reading those files recover as many identifiers as parsers that read everything, while the
verifier keeps every hallucinated LLM quote out of the evidence. Future work: train the planner for the top of the
list with a pairwise or listwise loss, since AP is the metric; use the fine-tuned LLM directly as the ranker and
inside the analyst agent, where its knowledge of which artifacts and tools examiners use may matter more than as a
prior; more labelled cases from public write-ups and examiner relevance judgments.

\begin{thebibliography}{99}

\bibitem{garfinkel2010} S. L. Garfinkel, ``Digital forensics research: The next 10 years,'' \emph{Digital
Investigation}, vol. 7, pp. S64--S73, 2010.

\bibitem{quick2014} D. Quick and K.-K. R. Choo, ``Impacts of increasing volume of digital forensic data: A survey
and future research challenges,'' \emph{Digital Investigation}, vol. 11, no. 4, pp. 273--294, 2014.

\bibitem{llava} H. Liu, C. Li, Q. Wu, and Y. J. Lee, ``Visual instruction tuning,'' in \emph{Advances in Neural
Information Processing Systems}, vol. 36, 2023, pp. 34\,892--34\,916.

\bibitem{qwen25vl} S. Bai \emph{et al.}, ``Qwen2.5-VL technical report,'' 2025.

\bibitem{kalai2025} A. T. Kalai, O. Nachum, S. S. Vempala, and E. Zhang, ``Why language models hallucinate,''
2025.

\bibitem{carrier2005} B. Carrier, \emph{File System Forensic Analysis}. Boston, MA: Addison-Wesley, 2005.

\bibitem{vaswani2017} A. Vaswani \emph{et al.}, ``Attention is all you need,'' in \emph{Advances in Neural
Information Processing Systems}, 2017, pp. 5998--6008.

\bibitem{qlora} T. Dettmers, A. Pagnoni, A. Holtzman, and L. Zettlemoyer, ``QLoRA: Efficient finetuning of
quantized LLMs,'' in \emph{Advances in Neural Information Processing Systems}, 2023.

\bibitem{hu2022lora} E. J. Hu \emph{et al.}, ``LoRA: Low-rank adaptation of large language models,'' in
\emph{International Conference on Learning Representations}, 2022.

\bibitem{ji2023} Z. Ji \emph{et al.}, ``Survey of hallucination in natural language generation,'' \emph{ACM
Computing Surveys}, vol. 55, no. 12, pp. 1--38, 2023.

\bibitem{wickramasekara2024} A. Wickramasekara and M. Scanlon, ``A framework for integrated digital forensic
investigation employing AutoGen AI agents,'' in \emph{2024 12th International Symposium on Digital Forensics and
Security (ISDFS)}. IEEE, 2024, pp. 1--6.

\bibitem{khatiwala2026} J. P. Khatiwala, D. K. N. Addai, and W. Xu, ``Evaluating the reliability of digital
forensic evidence discovered by large language model: A case study,'' \emph{arXiv preprint arXiv:2602.20202},
2026.

\bibitem{forensicllm} B. Sharma, J. Ghawaly, K. McCleary, A. M. Webb, and I. Baggili, ``ForensicLLM: A local large
language model for digital forensics,'' \emph{Forensic Science International: Digital Investigation}, vol. 52, p.
301872, 2025.

\bibitem{guwor2026} B. Guwor and M. Shabaz, ``Multi-modal deep learning for image forgery detection: A synergistic
fusion approach combining visual artifacts and metadata consistency analysis,'' \emph{Multimedia}, vol. 2, p. 11,
2026.

\bibitem{mittal2021} G. Mittal, P. Korus, and N. Memon, ``FiFTy: Large-scale file fragment type identification
using convolutional neural networks,'' \emph{IEEE Transactions on Information Forensics and Security}, vol. 16, pp.
28--41, 2021.

\bibitem{clip} A. Radford \emph{et al.}, ``Learning transferable visual models from natural language supervision,''
in \emph{Proceedings of the 38th International Conference on Machine Learning (ICML)}, 2021, pp. 8748--8763.

\bibitem{sbert} N. Reimers and I. Gurevych, ``Making monolingual sentence embeddings multilingual using knowledge
distillation,'' in \emph{Proceedings of the 2020 Conference on Empirical Methods in Natural Language Processing
(EMNLP)}, 2020, pp. 4512--4525.

\bibitem{oord2018} A. van den Oord, Y. Li, and O. Vinyals, ``Representation learning with contrastive predictive
coding,'' \emph{arXiv preprint arXiv:1807.03748}, 2018.

\bibitem{guo2017} C. Guo, G. Pleiss, Y. Sun, and K. Q. Weinberger, ``On calibration of modern neural networks,'' in
\emph{Proceedings of the 34th International Conference on Machine Learning (ICML)}, 2017, pp. 1321--1330.

\bibitem{wang2020} W. Wang, D. Tran, and M. Feiszli, ``What makes training multi-modal classification networks
hard?'' in \emph{Proceedings of the IEEE/CVF Conference on Computer Vision and Pattern Recognition (CVPR)}, 2020,
pp. 12695--12705.

\bibitem{cheng2016} H.-T. Cheng \emph{et al.}, ``Wide \& deep learning for recommender systems,'' in
\emph{Proceedings of the 1st Workshop on Deep Learning for Recommender Systems (DLRS)}, 2016, pp. 7--10.

\bibitem{bachlechner2021} T. Bachlechner, B. P. Majumder, H. Mao, G. Cottrell, and J. McAuley, ``ReZero is all you
need: Fast convergence at large depth,'' in \emph{Proceedings of the 37th Conference on Uncertainty in Artificial
Intelligence (UAI)}, 2021, pp. 1352--1361.

\bibitem{chawla2002} N. V. Chawla, K. W. Bowyer, L. O. Hall, and W. P. Kegelmeyer, ``SMOTE: Synthetic minority
over-sampling technique,'' \emph{Journal of Artificial Intelligence Research}, vol. 16, pp. 321--357, 2002.

\bibitem{zhang2018} H. Zhang, M. Cisse, Y. N. Dauphin, and D. Lopez-Paz, ``mixup: Beyond empirical risk
minimization,'' in \emph{International Conference on Learning Representations (ICLR)}, 2018.

\bibitem{xie2020} Q. Xie, M.-T. Luong, E. Hovy, and Q. V. Le, ``Self-training with noisy student improves ImageNet
classification,'' in \emph{Proceedings of the IEEE/CVF Conference on Computer Vision and Pattern Recognition
(CVPR)}, 2020, pp. 10687--10698.

\bibitem{ratcliff1988} J. W. Ratcliff and D. E. Metzener, ``Pattern matching: The gestalt approach,'' \emph{Dr.
Dobb's Journal}, vol. 13, no. 7, p. 46, 1988.

\end{thebibliography}

\end{document}
